\subsection{Bonus SVM}
\label{subsec:svm-bonus}

\todo{Hapus subbab bonus yang tidak dikerjakan.}

\subsubsection{Visualisasi Bidang Pemisah}
\todo{Gambar bidang pemisah, \textit{margin}, dan \textit{support vector} untuk tiap kernel pada dua fitur atau hasil reduksi ke 2D. Jelaskan fitur yang dipilih dan tampilkan tiga gambarnya.}

\begin{figure}[H]
	\centering
	\placeholdergambar{Visualisasi Bidang Pemisah}
	\caption{Visualisasi Bidang Pemisah SVM}
	\label{fig:svm-bonus-0}
\end{figure}

\subsubsection{Algoritma Optimasi di Luar Materi Kelas}
\todo{Jika ada: jelaskan algoritmanya, mengapa dapat mengoptimasi model, bukti perbaikannya (sebelum vs sesudah), dan \textbf{referensinya (tanpa referensi bonus dianulir)}. Misalnya SMO dengan pemilihan \textit{working set} orde kedua dan \textit{kernel cache} \cite{platt1998smo,chang2011libsvm}.}
